%% =====================================================================
%%  B5-T-07-04 -- what B5 contributes to "The Mirror"
%%  -------------------------------------------------------------------
%%  LAYER A: APPEND-ONLY. Written once, then left alone. Material found
%%  only here sits in the `unique` slot and is protected by rule R10.
%% =====================================================================
%% @kind: source
%% @id: B5-T-07-04
%% @book: B5
%% @topic: T-07-04
%% @locator: ch. 2, pp. 36--69
%% @status: distilled
%% @unique: yes
%% @integrated: yes
%% @updated: 2026-09-19

\dossier{B5}{T-07-04}{\bookshort{B5}}{ch. 2, pp. 36--69}

\begin{distilled}
  \item Mirroring is named as the art of imitation and linked to isopraxism: a neurobehaviour in which we copy each other to comfort each other, in speech patterns, body language, vocabulary, tempo and tone, mostly unconsciously, as a sign of bonding, being in sync, and establishing kinship.
  \item The technique: repeat the last one to three critical words the counterpart said, with a slightly upward tone, then be silent for several seconds and let the counterpart respond.
  \item The instruction is to fear silence: most talking is done by people afraid of quiet, and the void after a mirror is what pulls the elaboration out.
  \item The source pairs mirrors with the late-night FM DJ voice --- deep, soft, slow, reassuring --- as the delivery vehicle that keeps the mirror from reading as challenge, credited with de-escalating a confrontational call.
  \item A restaurant study is cited: waiters who mirrored their customers' orders earned materially larger tips than those who merely confirmed them, presented as evidence that imitation is read as rapport rather than as technique.
\end{distilled}

\begin{unique}
  \item The mirror-plus-silence pair stated as one tool: the mirror opens the valve and the silence does the extraction --- neither alone is the technique.
  \item The error taxonomy from the field: mirroring the wrong word (the incident instead of the stake) signals non-listening; stacking mirrors reads as mockery; both cost more than not mirroring.
  \item The delivery constraint: the same words in a confrontational voice fail; the mirror is only as safe as the voice carrying it.
\end{unique}
